\documentclass{article}
\usepackage[a4paper,margin=2cm]{geometry}
\usepackage{amsmath, amssymb, amsthm}
\usepackage{booktabs}
\usepackage{graphicx}
\newtheorem{proposition}{Proposition}
\theoremstyle{remark}
\newtheorem*{remark}{Remark}

\title{An explicit $(n,k)$ at which the optimal strategy is neither (A) nor (B)}
\author{}
\date{2026-04-16}

\begin{document}
\maketitle

\section*{Setting}

We consider the all-heads-wins game with $n\in\mathbb N$ coins, each landing
heads with probability $p\in(0,1)$ independently. In every round, all
remaining coins are tossed, and at least one coin must be set aside;
set-aside coins are never re-tossed. The game is won iff every set-aside
coin shows heads.

Write $w_{n,p}$ for the probability of winning with $n$ coins under
optimal play, and set $w_{0,p}:=1$. Denote the two standard strategies:
\begin{description}
  \item[(A)] set aside exactly one head, unless all coins are heads;
  \item[(B)] set aside every head (set aside one tails iff $k=0$).
\end{description}
If, in the current round, $n$ coins are tossed and $k\in\{1,\dots,n\}$ land
heads, a player may set aside any $i\in\{1,\dots,k\}$ heads (setting aside
a tails causes an immediate loss). After this round, $n-i$ coins remain and
the continuation value is $w_{n-i,p}$. Hence, by elementary conditioning,
\begin{equation}\label{eq:bellman}
  w_{n,p} \;=\; p^n + \sum_{j=1}^{n-1} \binom{n}{j} p^{\,n-j} (1-p)^j
  \max_{j\le m\le n-1} w_{m,p},
\end{equation}
where $j$ is the number of tails in the round.

The optimal number of heads to set aside at $(n,k)$ is therefore
\begin{equation}\label{eq:policy}
  i^{*}(n,k)\;\in\;\arg\max_{1\le i\le k} w_{n-i,p},
\end{equation}
and $i^{*}(n,k)=1$ iff strategy (A) is optimal at $(n,k)$, while
$i^{*}(n,k)=k$ iff strategy (B) is optimal at $(n,k)$.

\section*{Statement and proof}

\begin{proposition}
  Let $p=\tfrac{21}{50}$. At $(n,k)=(11,3)$ neither strategy {\rm (A)} nor
  strategy {\rm (B)} is optimal; the unique optimum is $i^{*}(11,3)=2$.
\end{proposition}

\begin{proof}
At $(n,k)=(11,3)$ the three available actions are $i\in\{1,2,3\}$ with
continuation values $w_{10,p}$, $w_{9,p}$, $w_{8,p}$. By \eqref{eq:policy}
it suffices to show
\begin{equation}\label{eq:need}
  w_{9,p} \;>\; w_{8,p} \qquad\text{and}\qquad w_{9,p} \;>\; w_{10,p}.
\end{equation}

Since $p=\tfrac{21}{50}\in\mathbb Q$, iterating \eqref{eq:bellman} in exact
rational arithmetic yields $w_{m,p}\in\mathbb Q_{\ge 0}$ for every
$m\in\mathbb N_0$. Starting from $w_{0,p}=1$, $w_{1,p}=\tfrac{21}{50}$, we
obtain the reduced representations (each step uses only the previously
computed values in \eqref{eq:bellman}; the suffix-max in each term is
determined by comparing the already-computed fractions):
\begin{align*}
  w_{1,p}  &= \tfrac{21}{50}, \qquad
  w_{2,p}   = \tfrac{11907}{31250}, \qquad
  w_{3,p}   = \tfrac{711920853}{1953125000}, \qquad
  w_{4,p}   = \tfrac{1093903231235157}{3051757812500000},\\[2pt]
  w_{5,p}  &= \tfrac{68010577069286354015457}{190734863281250000000000},\\[2pt]
  w_{6,p}  &= \tfrac{530643781539235033209174164045859}
                     {1490116119384765625000000000000000},\\[2pt]
  w_{7,p}  &= \tfrac{414482703926927350301614898070367347374127633}
                     {1164153218269348144531250000000000000000000000}.
\end{align*}
The three values relevant for \eqref{eq:need} are
\[
\resizebox{\linewidth}{!}{$\displaystyle
  w_{8,p}  = \frac{2023853327938433890470824968643571135424092441889732878537}
                   {5684341886080801486968994140625000000000000000000000000000},
$}
\]
\[
\resizebox{\linewidth}{!}{$\displaystyle
  w_{9,p}  = \frac{3952857618215383133100247584897156447999263365727606978962326704682729081}
                   {11102230246251565404236316680908203125000000000000000000000000000000000000},
$}
\]
\[
\resizebox{\linewidth}{!}{$\displaystyle
  w_{10,p} = \frac{77203229003644966332027370229551549794100203632172810638275046062922260090714793610061923}
                   {216840434497100886801490560173988342285156250000000000000000000000000000000000000000000000}.
$}
\]
Subtracting and reducing to lowest terms gives
\begin{equation}\label{eq:d98}
\resizebox{\linewidth}{!}{$\displaystyle
  w_{9,p}-w_{8,p}
    = \frac{19087085629440774417568015181574124082815161722450569748579682729081}
             {11102230246251565404236316680908203125000000000000000000000000000000000000}
$}
\end{equation}
and
\begin{equation}\label{eq:d910}
\resizebox{\linewidth}{!}{$\displaystyle
  w_{9,p}-w_{10,p}
    = \frac{1021352124235486336840412971037080885408979694513169582897387912292272566456389938077}
             {216840434497100886801490560173988342285156250000000000000000000000000000000000000000000000}.
$}
\end{equation}
Both fractions in \eqref{eq:d98} and \eqref{eq:d910} have a positive
integer numerator and a positive integer denominator, hence are strictly
positive rationals. This establishes \eqref{eq:need}, so
$i^{*}(11,3)=2$.
\end{proof}

\begin{remark}
  The value $w_{9,p}$ at $p=\tfrac{21}{50}$ is a strict local maximum of
  $m\mapsto w_{m,p}$: by \eqref{eq:d98}, $w_{9,p}>w_{8,p}$, and by
  \eqref{eq:d910}, $w_{9,p}>w_{10,p}$. The optimal action at $(11,3)$
  aims precisely at this peak---set aside $i=2$ heads so that $9$ coins
  remain in play, neither $n-1=10$ (strategy~A) nor $n-k=8$
  (strategy~B).
\end{remark}

\end{document}
